\documentclass[12pt,a4paper]{article}
\usepackage[brazilian]{babel}
\usepackage{amsmath,amssymb}
\usepackage{graphicx}
\usepackage[margin=2.5cm]{geometry}
\usepackage{hyperref}

\title{Um artigo de exemplo}
\author{Seu Nome}
\date{\today}

\begin{document}
\maketitle

\begin{abstract}
Este é um modelo de artigo com seções, equações, tabela e referências.
\end{abstract}

\section{Introdução}
\label{sec:intro}
O \LaTeX{} é ótimo para documentos científicos~\cite{lamport1994}.
Veja a Equação~\eqref{eq:euler} e a Seção~\ref{sec:resultados}.

\begin{equation}
  e^{i\pi} + 1 = 0
  \label{eq:euler}
\end{equation}

\section{Resultados}
\label{sec:resultados}

\begin{table}[h]
  \centering
  \begin{tabular}{lcr}
    \hline
    Item & Quantidade & Valor \\
    \hline
    A & 1 & 10,00 \\
    B & 2 & 20,00 \\
    \hline
  \end{tabular}
  \caption{Uma tabela simples.}
\end{table}

\[
  \int_0^\infty e^{-x^2}\,dx = \frac{\sqrt{\pi}}{2}
\]

\bibliographystyle{plain}
\bibliography{referencias}

\end{document}
